\documentclass[book.tex]{subfiles}

\begin{document}

\chapter{Perturbation Theory}

\label{chap:perturbation_theory}
\noindent\rule{11cm}{0.4pt} \\
\textbf{Prerequisites:} \autoref{chap:schrodinger} \\
\textbf{Difficulty Level:} ** \\
\noindent\rule{11cm}{0.4pt}
\vspace{5mm}

This chapter provides an introduce to perturbation theory. Broadly stated, perturbation theory is a technique for approximating hard to compute quantities by adding on additional terms (perturbations) to simpler physical theories. Perturbation theory is related to classical approximation techniques such as Taylor series expansions, but with the crucial difference that perturbation theory techniques may sometimes be fruitfully used for divergent series.

\section{Time Independent Perturbation Theory}

Suppose that we have a time independent Hamiltonian $H$ that can be written as
\begin{align}
    H &= H_0 + \epsilon V
\end{align}
Assume that for $H_0$ we have solved for the eigenvalues and eigenstates of the operator (for example, $H_0$ may be a simple system like the Harmonic oscillator.) We then seek solutions to $H$ starting from solutions to $H_0$ and expanding in powers of $\epsilon$

%\section{Time Dependent Perturbation Theory}

\end{document}